\begin{lemma}
\label{lemma-support-base-change}
Let $R \to R'$ be a ring map and let $M$ be a finite $R$-module.
Then $\text{Supp}(M \otimes_R R')$ is the inverse image of
$\text{Supp}(M)$.
\end{lemma}

\begin{proof}
Write the original map as $\varphi:R\to R'$ and the induced
spectrum map as $f:\Spec(R')\to\Spec(R)$.
Fix an original prime $\mathfrak p'\subset R'$ and its
contraction $\mathfrak p=\varphi^{-1}(\mathfrak p')$.
The local map is
$$
R_{\mathfrak p}\longrightarrow R'_{\mathfrak p'},
\qquad r/t\longmapsto\varphi(r)/\varphi(t).
$$
All denominators are valid, because $t\notin\mathfrak p$
implies $\varphi(t)\notin\mathfrak p'$. Its inverse image
of $\mathfrak p'R'_{\mathfrak p'}$ is
$\mathfrak pR_{\mathfrak p}$: a fraction in the former
ideal has numerator in $\mathfrak p'$, and contraction
of that numerator is the stated condition in $R$.

For any original module $M$, not yet assumed finite, the
exact localization-tensor comparison is
$$
(M\otimes_RR')_{\mathfrak p'}
\longrightarrow M_{\mathfrak p}\otimes_{R_{\mathfrak p}}R'_{\mathfrak p'},
\qquad \frac{m\otimes s}{u}\longmapsto\frac m1\otimes\frac su.
$$
Its inverse is
$$
\frac mr\otimes\frac su
\longmapsto\frac{m\otimes s}{\varphi(r)u},
\qquad r\notin\mathfrak p,\quad u\notin\mathfrak p'.
$$
The forward map is induced by the balanced map to the
right-hand module and then by its invertible action of
every $u\notin\mathfrak p'$. For the inverse, if
$v(r'm-rm')=0$ with $v\notin\mathfrak p$, multiplying
the corresponding numerator difference by
$\varphi(v)\notin\mathfrak p'$ gives zero. This checks
the equivalence relation in $M_{\mathfrak p}$. If
$w(u's-us')=0$ with $w\notin\mathfrak p'$, the same
$w$ kills the tensor numerator difference, checking the
other localization relation. Additivity uses common
denominators and tensor additivity. For $a/t\in R_{\mathfrak p}$,
the two balanced expressions map respectively to
$$
\frac{am\otimes s}{\varphi(t)\varphi(r)u}
=\frac{m\otimes\varphi(a)s}{\varphi(r)\varphi(t)u}.
$$
This proves balancing with every original factor. Composing
in one direction gives the original fraction immediately;
in the other, $m/1\otimes s/(\varphi(r)u)=m/r\otimes s/u$.
Thus the maps are inverse. In particular, for every $M$
and every $R\to R'$,
$$
\operatorname{Supp}_{R'}(M\otimes_RR')
\subseteq f^{-1}\operatorname{Supp}_R(M).
$$

The original residue-field map is the injective map
$$
\iota:\kappa(\mathfrak p)\longrightarrow\kappa(\mathfrak p'),
\qquad
\frac{r+\mathfrak p}{t+\mathfrak p}
\longmapsto
\frac{\varphi(r)+\mathfrak p'}{\varphi(t)+\mathfrak p'}.
$$
It exists and is injective because $R/\mathfrak p\to
R'/\mathfrak p'$ is injective and every displayed
nonzero denominator stays nonzero. There are exact maps
$$
\begin{split}
(M\otimes_RR')\otimes_{R'}\kappa(\mathfrak p')
&\longrightarrow M\otimes_R\kappa(\mathfrak p'),\\
(m\otimes s)\otimes\lambda
&\longmapsto m\otimes\overline s\lambda,
\end{split}
$$
with inverse $m\otimes\lambda\mapsto(m\otimes1)\otimes\lambda$,
and
$$
\begin{split}
(M\otimes_R\kappa(\mathfrak p))
\otimes_{\kappa(\mathfrak p)}\kappa(\mathfrak p')
&\longrightarrow M\otimes_R\kappa(\mathfrak p'),\\
(m\otimes\alpha)\otimes\beta
&\longmapsto m\otimes\iota(\alpha)\beta,
\end{split}
$$
with inverse $m\otimes\beta\mapsto(m\otimes1)\otimes\beta$.
The tensor balancing relations show independence and give
both composites on the displayed generators. Extension
along the field inclusion detects a nonzero vector space:
if $B$ is a basis, the map from the direct sum of copies
of $\kappa(\mathfrak p')$ indexed by $B$, sending its
$b$-coordinate $\lambda$ to $b\otimes\lambda$, is an
isomorphism, with inverse given by the finite basis
expansion of each vector. Thus it is zero exactly when
$B$ is empty.

For completeness, for any local ring $(A,\mathfrak m)$
and module $N$, the map
$N\otimes_A A/\mathfrak m\to N/\mathfrak mN$,
$n\otimes\overline a\mapsto an+\mathfrak mN$, has inverse
$n+\mathfrak mN\mapsto n\otimes1$; balancing and the
relations $\mathfrak mN=0$ in both targets prove the claim.
Applying this after the localization comparisons gives the
original fibre quotient
$$
M\otimes_R\kappa(\mathfrak p)
\simeq M_{\mathfrak p}/\mathfrak pM_{\mathfrak p},
$$
and the corresponding quotient for $M\otimes_RR'$ at
$\mathfrak p'$. Explicitly,
$m\otimes(a/t+\mathfrak pR_{\mathfrak p})$ maps to
$am/t+\mathfrak pM_{\mathfrak p}$, with inverse
$m/t+\mathfrak pM_{\mathfrak p}\mapsto
m\otimes(1/t+\mathfrak pR_{\mathfrak p})$.

Now restore the original finiteness hypothesis. The module
$M_{\mathfrak p}$ is finite over the local ring, so
Nakayama's Lemma \ref{lemma-NAK} says that it is nonzero
exactly when this fibre quotient is nonzero. If
$\mathfrak p\in\operatorname{Supp}_R(M)$, the field
extension and the displayed fibre maps show that the
original quotient
$$
(M\otimes_RR')_{\mathfrak p'}/
\mathfrak p'(M\otimes_RR')_{\mathfrak p'}
$$
is nonzero. Therefore its module is nonzero. Together
with the already proved reverse inclusion, this proves
the stated lemma, keeping both original primes and all maps.

\medskip\noindent
There are two further exact forms. First, for an arbitrary
$M$ put
$$
F_R(M)=\{\mathfrak p\in\Spec(R):
M\otimes_R\kappa(\mathfrak p)\ne0\}.
$$
The field-extension argument above did not use finiteness
or flatness. Hence for every ring map it proves
$$
F_{R'}(M\otimes_RR')=f^{-1}F_R(M).
$$
The localization comparison always gives
$F_R(M)\subseteq\operatorname{Supp}_R(M)$.
Consequently localization support commutes with every
base change if and only if
$$
F_R(M)=\operatorname{Supp}_R(M).
$$
For sufficiency, a prime in the inverse image of support
is in the inverse image of $F_R(M)$, hence in the new
fibre set and therefore in the new support. The general
reverse inclusion was proved above. For necessity, apply
the asserted support equality to the particular original
map $R\to\kappa(\mathfrak p)$. The field's unique prime
contracts to $\mathfrak p$, and its module is nonzero
exactly when $M\otimes_R\kappa(\mathfrak p)\ne0$.
This proves both directions of the criterion.

This criterion strictly extends finite generation. An
arbitrary direct sum of finite modules satisfies it.
Localization commutes with direct sums by the maps
$(m_i)_i/s\mapsto(m_i/s)_i$ and their inverse obtained
from the finite sum of the localized inclusions.
For coordinates $m_i/s_i$ indexed by a finite set $E$,
the inverse has denominator $\prod_{i\in E}s_i$ and
$i$-th numerator $m_i\prod_{j\in E,\,j\ne i}s_j$.
Tensor products commute with these direct sums by
$(m_i)_i\otimes c\mapsto(m_i\otimes c)_i$, with inverse
given by the sum of the tensor products of the original
inclusions. Balancing verifies these maps, and both
composites are identities on elements of finite support.
A direct sum is nonzero exactly when one summand is;
Nakayama for each finite summand therefore proves the
criterion for the sum. For a nonzero ring, a free module
on an infinite set is not finite: any finite collection
of vectors involves only finitely many basis coordinates,
so cannot generate a basis vector outside their union.
It supplies the strictness example.

\medskip\noindent
Second, if the original map $R\to R'$ is flat, support
commutes with this base change for every $M$. The local
map $R_{\mathfrak p}\to R'_{\mathfrak p'}$ is flat:
tensor with $R'$ preserves injections, and subsequent
localization at $\mathfrak p'$ preserves them. For any
$R_{\mathfrak p}$-module $N$, the comparison
$N\otimes_RR'_{\mathfrak p'}\simeq
N\otimes_{R_{\mathfrak p}}R'_{\mathfrak p'}$
sends each original elementary tensor to itself.
Its inverse is balanced over $R_{\mathfrak p}$ since
every $t\notin\mathfrak p$ acts invertibly in both factors;
thus the comparison follows with no injectivity assumption
on the ring map.

A flat local ring map $(A,\mathfrak m)\to(B,\mathfrak n)$
detects nonzero arbitrary modules. Choose $x\ne0$ in such
a module $N$ and put $J=\operatorname{Ann}_A(x)$.
The map $A/J\hookrightarrow N$, $a+J\mapsto ax$, is
injective. Flatness gives $B/JB\hookrightarrow N\otimes_AB$.
Here $J\subseteq\mathfrak m$ and
$JB\subseteq\mathfrak mB\subseteq\mathfrak n$, so
$B/JB\ne0$. This proves faithful flatness of the local
map, and applies to the original localization map above.
Its tensor comparison proves
$$
M_{\mathfrak p}\ne0
\quad\Longleftrightarrow\quad
(M\otimes_RR')_{\mathfrak p'}\ne0
$$
for every $M$, proving the asserted flat-base-change extension.

Flatness of the module alone does not imply the universal
support criterion. Retain the example
$R=\mathbf Z$, $M=\mathbf Q$, $R'=\mathbf F_p$ for a
prime integer $p$, with the quotient map. The module
$\mathbf Q$ is flat by its exact localization comparison,
proved in Example \ref{example-flat-infinite-intersections}.
At every prime of $\mathbf Z$, its localization is still
$\mathbf Q$: the map sending an original fraction $q/s$
to the same rational number is inverse to $q\mapsto q/1$.
Thus its localization support is all of $\Spec(\mathbf Z)$.
But
$$
q\otimes\overline a=(q/p)\otimes p\overline a=0
$$
in $\mathbf Q\otimes_{\mathbf Z}\mathbf F_p$, so the
new support is empty, whereas the inverse image of the
old support is the nonempty spectrum of $\mathbf F_p$.
The fibre set is exactly $\{(0)\}$: all nonzero prime
fibres vanish as just proved, while multiplication gives
$\mathbf Q\otimes_{\mathbf Z}\mathbf Q\simeq\mathbf Q$,
with inverse $q\mapsto q\otimes1$. For a fraction $a/b$,
$q\otimes a/b=(qa/b)\otimes1$ because $b$ is invertible
on both sides of the tensor, proving the inverse formula.
Moreover the annihilator of $\mathbf Q$ is detected by
its element $1$. Thus finite annihilator detection does
not imply the universal support criterion. Empty spectra
and zero modules satisfy all the stated identities as well.
\end{proof}